% -*- TeX-engine: luatex; -*-
\documentclass{article}
\usepackage{ttn}
\usepackage{booktabs}
\usepackage{amsmath}
\usepackage{siunitx}
% \usepackage{mathtools}
\usepackage{colonequals} % because unicode-math has broken "::=" somehow
%
% \usepackage[raggedrightboxes]{ragged2e}
\usepackage[Ragged, size=footnote, shape=up, alerton]{sidenotesplus}
\usepackage{subcaption}
\DeclareCaptionStyle{sidecaption}{font=footnotesize, justification=raggedright}
\usepackage{graphicx}
\usepackage[dvipsnames]{xcolor}
\usepackage[colorlinks=true,linkcolor=Blue]{hyperref}
%
%% \defn takes an optional star and an optional argument.
%% If the optional argument is present, it is used for the marginal note; if the argument is missing, the default is the mandatory argument. if the command is starred, no marginal note is made. 
\NewDocumentCommand\defn{sO{#3}m}{\textit{#3}\IfBooleanF{#1}{\sidenote*{\textcolor{MidnightBlue}{#2}}}}
%\newcommand{\defn}[2][#2]{\textit{#2}\sidenote*{#1}}
\newcommand{\isdef}{\stackrel{\text{def}}{=}}
\newcommand{\set}[1]{\mathbf{#1}}
\newcommand{\R}{\set{R}}
\newcommand{\V}{\mathcal{V}}
\newcommand{\W}{\mathcal{W}}
\newcommand{\vv}[1]{\symbf{#1}}
%% 
\title{Automatic Differentiation}
\author{James Geddes}
\date{\today}
%%
\begin{document}%\maketitle
\maketitle

\section{The gradient}
\sidenote*{\emph{Notation:} In this section, a boldface symbol will indicate an
element of a vector space. For example, $\vv{v}\in\mathcal{V}$ indicates that
$\vv{v}$ is an element of the vector space~$\mathcal{V}$.  Note that this
notation does not distinguish between different vector spaces even
though an element of one vector space is not the same thing as an
element of a different vector space.}
\sidenote*{A real vector space is a set that comes with a
  notion of ``addition'' and ``multiplication by real numbers.''  A
  \emph{vector} is an element of a vector space.  Sometimes people use
  the word ``vector'' to mean ``a tuple of real numbers'' but that is
  not correct.}\sidenote*{More carefully, a \emph{real vector space}
  is a set, $\V$, together with two operations:
  ${+}:\V \times \V \to \V$ (“addition”) and
  ${\cdot}:\R \times \V \to \V$ (“multiplication by a scalar”), and a
  distinguished element $ \vv{0}\in \V$, all of which satisfy the
  following rules:
\begin{enumerate}
\item
  For every $\vv{u}, \vv{v}, \vv{w} \in \V$,
  \begin{equation*}
    \begin{gathered}
      \vv{v} + \vv{w} = \vv{w} + \vv{v}, \\
      (\vv{u} + \vv{v}) + \vv{w} = \vv{u} + (\vv{v} + \vv{w}), \\
      \vv{0} + \vv{v} = \vv{v} + \vv{0} = \vv{v};
      \end{gathered}
  \end{equation*}
\item For every $\vv{v} \in \V$ there must exist an element $-\vv{v}$ such that
  \begin{equation*}
    \vv{v} + (-\vv{v}) = \vv{0}; 
  \end{equation*}
\item For every $\alpha, \beta \in \R$ and $\vv{v}, \vv{w} \in \V$,
  \begin{equation*}
    \begin{gathered}
      1 \cdot \vv{v} = \vv{v} \\
      \alpha \cdot (\beta \cdot \vv{v}) = (\alpha \beta) \cdot \vv{v}, \\
      (\alpha + \beta) \cdot \vv{v} = \alpha \cdot \vv{v} + \beta \cdot \vv{v}, \\
      \alpha \cdot (\vv{v} + \vv{w}) = \alpha \cdot \vv{v} + \alpha \cdot \vv{w}.
    \end{gathered}
  \end{equation*}
\end{enumerate}
} Suppose that $f(x)$ is a real-valued function on the reals. The
ordinary derivative of $f$ at $x$ is:
\begin{equation}\label{eq:deriv}
  \frac{df}{dx} = \lim_{\delta\to0}\frac{f(x + \delta)-f(x)}{\delta}.
\end{equation}
I am sure you are aware of the derivatives of elementary functions
as well as properties of the derivative such as the chain rule.  Here
is one simple property. Suppose the Taylor expansion of $f$ in the
vicinity of $x=x_0$ is
\begin{equation*}
  f(x) = a + b (x-x_0) + \frac{1}{2}c(x-x_0)^2 + \dotsb. 
\end{equation*}
Then the derivative of $f$ at $x_0$ is just~$b$, the coefficient of
the linear term in the expansion.

The idea of the gradient is to generalise this ordinary derivative.
Instead of functions from $\R$ to $\R$, we are interested in functions
from one vector space to another.  Can we generalise the derivative to this
setting?

Fix, once and for all, two finite-dimensional, real vector spaces,
$\mathcal{V}$ and $\mathcal{W}$, and let
$f\colon \mathcal{V}\to\mathcal{W}$ be a smooth map from $\mathcal{V}$
to~$\mathcal{W}$.\sidenote{By ``smooth'' we mean that we shall assume that we
  can take derivatives whenever we like.}  We shall construct the
\emph{gradient} of $f$, which will be the generalisation of the
ordinary derivative.

\subsection{Two special case}

\emph{Example:} Suppose we simply things slightly by choosing
$\V=\set{R}$. Just for this example, I will write the map, $\vv{f}$,
in boldface, so that $\vv{f}\colon\R\to\W$ is a map from $\R$
to~$\W$. Thus, $\vv{f}$ describes a path in
$\W$ parameterised by a number, $t\in \R$.  One might think of
$\vv{f}(t)$ as ``a position in $\W$ at time~$t$.''  The derivative of
this map is defined by the same formula as in eq.~\eqref{eq:deriv}; that
is:
\begin{equation*}
  \frac{d\vv{f}}{dt} = \lim_{\delta\to0}\frac{\vv{f}(t + \delta)-\vv{f}(t)}{\delta}.
\end{equation*}
Notice that because $\mathcal{W}$ is a vector space, the expression on the
right-hand side is well-defined. That is to say, it is permissible to
write $\vv{f}(t + \delta)-\vv{f}(x)$ because the notion of the subtraction
of one vector from another is part of the required machinery of a
vector space.  Likewise, the division by $\delta$, a real number, is also
an operation available in~$\W$. The result of the right-hand side is a
vector in~$\W$ and so the derivative is, for each $t\in\R$, a vector
in~$\W$.

Sometimes one writes $\vv{f}'$ (``$f$ prime'') for the derivative
of~$\vv{f}$.\sidenote{An alternative notation, especially when the
  parameter is thought of as ``time,'' is $\dot{\vv{f}}$ (``$f$
  dot'').} Since $\vv{f}'$ ``lives in'' $\W$, we can also imagine
that, to each point on the path $\vv{f}(t)$ traced out in $\W$, we assign
the vector~$\vv{f}'(t)$ (see figure~\ref{fig:velocity}).  If
$\vv{f}(t)$ is a ``position at time $t$'', then $\vv{f}'(t)$ might be
thought of as ``the velocity at time~$t$.''

\emph{Example:} Suppose instead that $\V$ is an arbitrary vector space
but now $\W = \set{R}$.  Now $f$ is a real-valued function on~$\V$; that
is to say, $f(\vv{x})$ is a real number for each $\vv{x}\in\V$.  (In a
sense, this is the opposite of the previous example.)  On the face of
it, one might imagine that we can directly adapt eq.~\eqref{eq:deriv}
to this circumstance as well.  After all, there should be no trouble
in computing, say, the subtraction on the right-hand side of
eq.~\eqref{eq:deriv} since both terms are simply numbers.
\begin{marginfigure}
  \margincaption{An example path, $f(t)$ (in blue), and samples of its derivative
    (shown as red arrows).\label{fig:velocity}} \input{deriv-of-path.tex}
\end{marginfigure}
However, if one attempts actually to write down the formula, one arrives
at something like the following:
\begin{equation*}
  \frac{df}{d\vv{x}} = \lim_{\vv{\delta}\to0}\frac{f(\vv{x} +
    \vv{\delta})-f(\vv{x})}{\vv{\delta}}  \qquad\text{--- Not well formed!}
\end{equation*}
One immediate problem problem with this formula is that, on the
right-hand side, we are attempting to divide by a vector,
$\vv{\delta}$, and that operation is certainly not well-defined.  A
more subtle problem is that, since $\vv{x}$ is a vector, the notion of
a ``small change in $\vv{x}$'' must also be a vector; that is, the
derivative of $f$ now depends on the direction in which one is taking
the derivative.

It turns out that, if we are willing, for now, to go along with the
idea of a ``directional derivative,'' then it is possible to produce a
well-defined expression for such a thing. Fix a vector,
$\vv{\xi}$, and let $\delta>0$ be a number. Now consider the
expression:
\begin{equation*}
  \lim_{\delta\to0}\frac{f(\vv{x} +
    \delta\vv{\xi})-f(\vv{x})}{\delta}.
\end{equation*}
\emph{This} expression is well-defined: on the right, the denominator,
$\delta$, is simply a number; and the limit is taken as $\delta$ approaches
zero.  We call this limit the ``directional derivative of $f$
along $\vv{\xi}$'' and denote it by
$\nabla_{\vv{\xi}}f$.\sidenote{Other notations for the directional
  derivative are $D_{\vv{\xi}}f$, $\pounds_{\vv{\xi}}f$, and
  $\partial_{\vv{\xi}}f$.} An equivalent way to formulate the directional
derivative is
\begin{equation}
  \label{eq:grad}
  \nabla_{\vv{\xi}}f \isdef \left.\frac{d}{d\lambda} f(\vv{x}+\lambda\vv{\xi})\right|_{\lambda=0}.
\end{equation}

The directional derivative, $\nabla_{\vv{\xi}}f$, of $f$, at $\vv{x}$,
is a number.  This number depends on the direction,~$\vv{\xi}$.  You
might expect that the number would depend in some complicated way on
the direction.  Remarkably, however, the derivative is \emph{linear}
in~$\vv{\xi}$.\sidenote{To show that the directional derivative is
  linear in the direction, one must show that
  $\nabla_{\alpha\vv{\xi}+\beta\vv{\zeta}}f = \alpha\nabla_{\vv{\xi}} f+ \beta\nabla_{\vv{\zeta}} f$
  for any two numbers $\alpha,\beta\in\set{R}$ and vectors $\vv{\xi},\vv{\zeta} \in\V$.
  %
  \sidepar
  To see this, let $s(\lambda)$ and $t(\lambda)$ be any two functions of
  $\lambda$ and consider the derivative, with respect to $\lambda$, of the
  function $f(\vv{x}+s(\lambda)\vv{\xi} + t(\lambda)\vv{\zeta})$. By the chain rule, one
  has:
\begin{equation*}
  \frac{d}{d\lambda}f(\vv{x}+s(\lambda)\vv{\xi} + t(\lambda)\vv{\zeta}) = \frac{\partial f}{\partial s} \frac{ds}{d\lambda} +
  \frac{\partial f}{\partial t} \frac{dt}{d\lambda}.
\end{equation*}
Now choose the special case of $s(\lambda)=\alpha\lambda$ and $t(\lambda)=\beta\lambda$, whence
\begin{equation*}
  \frac{d}{d\lambda}f(\vv{x}+\lambda(\alpha\vv{\xi} + \beta\vv{\zeta})) = \frac{\partial f}{\partial s} \alpha +
  \frac{\partial f}{\partial t} \beta.
\end{equation*}
Evaluating this at $\lambda=0$ one has, on the left, the directional
derivative along $\alpha\vv{\xi} + \beta\vv{\zeta}$; and, on the
right, the required linear combination of the directional derivatives
along $\vv{\xi}$ and~$\vv{\zeta}$.  } That is, there is, at each point
in $\V$, some linear map, $\vv{G}\colon\V\to\R$, such that, for
\emph{any} vector $\vv{\xi}$, we have
$\nabla_{\vv{\xi}}f = \vv{G}(\vv{\xi})$.  This $\vv{G}$ is called the
\emph{gradient} of $f$, written~$\nabla f$. 

Note that in some sense the gradient has two arguments: it depends on
$\vv{x}$, the point in~$\V$, and, for each $\vv{x}\in\V$, the gradient
at that point is a linear map, which ``consumes'' a value
$\vv{\xi}\in\V$ to produce a value in~$\R$.  (The dependence on
$\vv{x}$ is often omitted and the dependence on $\vv{\xi}$ is sometimes
shown as a ``dot product.'')

In summary: given a function $f(\vv{x})$ there is, at each point
$\vv{x}\in\V$, a linear map, $\nabla f\colon\V\to\R$, such that for any vector
$\vv{\xi}$,
\begin{equation*}
  \nabla_{\vv{\xi}} f = (\nabla f)(\vv{\xi}),
\end{equation*}
where on the right hand side, we mean to show the action of $\nabla f$ as a
linear map and we have suppresed the dependence on~$\vv{x}$. .

It is worth stepping back briefly and recalling the common story
about the method of gradient descent.  To find the minimum of a
function by the method of gradient descent, one computes the gradient
of the function, which is supposed to ``point uphill.''  That is, the
gradient is thought of as a vector in~$\V$.  But we have just
concluded that the gradient, $\nabla f$, is \emph{not} an element
of~$\V$; it is, instead, an element of $\mathcal{L}(\V, \R)$, the space of
linear maps from $\V$ to~$\R$.  Those are not the same thing.  What is
going on?

The space $\mathcal{L}(\V,\R)$ is itself a vector space.  This vector space is
called the \emph{dual} of~$\V$ and often denoted~$\V^*$.
A notable fact about $\V^*$ is that it has the same dimension as~$\V$.
Thus, if one happens to represent an element of an $n$-dimensional
space $\V$ by a tuple of $n$ numbers, then one might also represent an
element of $\V^*$ by a tuple of $n$ numbers.\sidenote{People often visually distinguish the two
  kinds of tuple by writing an element of $\V$ vertically--a ``column
  vector''---and writing an element of $\V^*$ horizontally--a ``row
  vector.'' Other terminology calls an element of $\V$ a
  \emph{contravariant vector} and an element of $\V^*$ a
  \emph{covariant vector}.}

If one squints, those look like the same thing.  Perhaps one can
obtain a ``real'' vector in $\V$ from a dual vector in $\V^*$ by
taking the components of the dual vector and simply pretending they
are the components of a vector?

Well, yes, you can, but the resulting vector will depend on the
specific basis chosen.  There is no ``natural'' way to
convert an element of $\V^*$ into an element of $\V$: some additional
structure is required, such as a choice of basis (or an inner
product).  Absent such additional structure, the conversion will
inevitably involve an arbitrary choice, and one ought to feel
uncomfortable that the result depends on that arbitrary choice.
 
Summarising the above examples, we have the following:
\begin{enumerate}
\item The derivative of a function $f\colon\R\to\R$ is, at each point, a
  number; that is, an element of $\R$.
\item The derivative of a function $f\colon \R\to\W$ is, at each point,
  a vector in~$\W$.
\item The derivative of a function $f\colon\V\to\R$ is, at each point, a
  dual vector in~$\V^*$; that is, a linear map $\V\to\R$.
\end{enumerate}

\subsection{The general case}

We return now to the general case of arbitrary, finite-dimensional
vector spaces $\V$ and~$\W$, and a (vector-valued) function $\vv{f}\colon\V\to\W$.  The
remarks above give us the notion of a $\W$-valued directional
derivative of~$\vv{f}$, where the direction is taken in~$\V$.  Then a
similar argument to one above shows that the directional derivative is
linear in the direction and hence arises from a linear map. In other
words, for a function,
\begin{equation*}
  \vv{f}\colon\V\to\W,
\end{equation*}
we obtain a gradient,
\begin{equation*}
  \nabla \vv{f}\colon \V\to\mathcal{L}(\V,\W),
\end{equation*}
by
\begin{equation}\label{eq:gradient}
  \nabla \vv{f} \colon \vv{x} \mapsto \vv{\xi} \mapsto \left.\frac{d}{d\lambda}\vv{f}(\vv{x}+\lambda \vv{\xi})\right|_{\lambda=0}  .
\end{equation}

\emph{Example:} Suppose $\vv{f}$ is a constant function, say $\vv{f}(\vv{x})=\vv{c}$,
where $\vv{c}\in\W$, for any~$\vv{x}\in\V$.  Then $\nabla \vv{f}$ is, at each point $\vv{x}\in\V$,
the zero map, $\nabla \vv{f}\colon \vv{v} \mapsto \vv{0}$.

\emph{Example:} Suppose $\vv{f}$ is itself a linear map, that is
$\vv{f}\in\mathcal{L}(\V,\W)$.  Consider the derivative on the right-hand side of
eq.~\eqref{eq:gradient} at some point $\vv{x}\in\V$ and given a direction
$\vv{\xi}\in\V$.  Since $\vv{f}$ is linear, we have
\begin{equation*}
  \begin{aligned}
  \left.\frac{d}{d\lambda}\vv{f}(\vv{x}+\lambda\vv{\xi})\right|_{\lambda=0}
    &=   \left.\frac{d}{d\lambda} \Bigl[\vv{f}(\vv{x}) +\lambda f(\vv{\xi})\Bigr] \right|_{\lambda=0} \\
    &= \vv{f}(\vv{\xi}).
  \end{aligned}
\end{equation*}
Since there is no dependence on $\vv{x}$, the gradient of a linear
function is constant; and that constant is the constant linear map
$\vv{f}$ itself.

\emph{Example:} Suppose that $\vv{v}_{(1)}, \dotsc, \vv{v}_{(n)}$ is a basis
for~$\V$, and $\vv{w}_{(1)}, \dotsc, \vv{w}_{(m)}$ is a basis for~$\W$. (I mean
that, $\vv{v}_{(1)}$, for example is the first basis vector in $\V$,
\emph{not} that it is the first component of some vector.)

Let $f^j$ (for $j=1,\dotsc,m$) be the components of $\vv{f}$ in the
basis of~$\W$, so that
$\vv{f} = \sum_j f^j \vv{w}_{(j)}$.\sidenote{You will have noted a
  convention in the placement of the indices. The basis vectors are
  enumerated by a ``down'' index in parentheses, such
  as~$\vv{w}_{(i)}$, whereas the components are indicated by an ``up''
  index, such as~$f^i$.  Contrariwise, I will denote a particular dual
  basis vector with an up index in parentheses, and the components of
  a dual vector by a down index.  The net result will be that
  summations will only ever occur over an index where one occurence is
  up and one is down.} What are the components of $\nabla f$ expressed in
these bases?

What does it actually mean to ``express $\nabla f$ in these bases''?
$\nabla f$~lives in $\mathcal{L}(\V,\W)$ and so we will need a basis
for~$\mathcal{L}(\V,\W)$.  What does a general element of
$\mathcal{L}(\V,\W)$ look like?  Pick some $\tilde{\vv{v}}\in\V^*$, the dual
of~$\V$, and some $\vv{w}\in\W$, and consider the operation, on an
element $\vv{x}\in\V$, that is ``act on $\vv{x}$ with $\tilde{\vv{v}}$
to produce a number, then multiply $\vv{w}$ by this number.''  This
operation is a linear map from $\V$ to~$\W$.  We write it as
$\tilde{\vv{v}}\otimes \vv{w}$, and the description of this map above is
equivalent to
\begin{equation*}
  (\tilde{\vv{v}}\otimes \vv{w})(\vv{x}) \isdef (\tilde{\vv{v}}(\vv{x})) \vv{w}.
\end{equation*}

Now suppose that $\tilde{\vv{v}}^{(1)}, \dotsc, \tilde{\vv{v}}^{(n)}$
is a basis for~$\V^*$ (the next paragraph will describe how to obtain
this basis).  We claim (without proof) that a basis for
$\mathcal{L}(\V,\W)$ is then given by the $n\times m$ objects
$\tilde{\vv{v}}^{(1)}\otimes \vv{w}_{(1)}, \tilde{\vv{v}}^{(2)}\otimes
\vv{w}_{(1)},\dotsc,\tilde{\vv{v}}^{(1)}\otimes \vv{w}_{(2)},\dotsc,
\tilde{\vv{v}}^{(n)}\otimes \vv{w}_{(m)}$.

The natural basis to choose for $\V^*$ is the \emph{dual basis} (dual
to the basis $\vv{v}_{(i)}$ of~$\V$) defined by the relations:
\begin{equation*}
  \tilde{\vv{v}}^{(i)} (\vv{v}_{(j)}) = \begin{cases}
    1\quad\text{if $i=j$,} \\
    0\quad\text{otherwise.}
  \end{cases}
\end{equation*}
In this basis, the components of any $\tilde{\vv{x}}\in\V^*$ are
straightforward to compute: they are $\tilde{x}_j =
\tilde{\vv{x}}(\vv{v}_{(j)})$. 

So then, since the $\vv{w}_{(j)}$ are constant, we have\sidenote{I am
  relying on the following lemma which I have not shown: for $f$ a
  real-valued function and $\vv{w}\in\W$ constant, we have
  $\nabla (f\vv{w}) = (\nabla f) \otimes \vv{w}$.}
\begin{equation*}
  \begin{aligned}
    \nabla f
    &= \nabla \bigl(\sum_j f^j \vv{w}_{(j)}\bigr) \\
    &= \sum_j (\nabla f^j) \otimes \vv{w}_{(j)}).
  \end{aligned}
\end{equation*}
Now, since $f^j$ is a real-valued function, $\nabla f^j$ is an element
of~$\V^*$, which we may write as  
\begin{equation*}
  \nabla f^j = \sum_i f^j{}_i \tilde{\vv{v}}^{(i)}.
\end{equation*}
That is to say, for fixed $j$, the value $f^j{}_i$ is the $i$th
component of $\nabla f^j$ in the dual basis. To isolate a single component,
we apply $\vv{v}_{(k)}$ to both sides.  On the right-hand side, we obtain
just $f^j{}_k$. Thus,
\begin{equation*}
  \begin{aligned}
    f^j{}_k &= (\nabla f^j)(\vv{v}_{(k)}) \\
    &= \nabla_{\vv{v}_{(k)}} f^j \\
    &= \frac{d}{d\lambda} f^j(x^1, \dotsc, x^k + \lambda, \dotsc, x^n)\Bigr|_{\lambda=0} \\
    &= \frac{\partial f^j}{\partial x^k}.
  \end{aligned}
\end{equation*}
(The second step follows from the definition of $\nabla$ in terms of the
directional derivative and the third step follows from the definition
of the directional derivative.)

The numbers $f^j{}_k$ are called the \emph{Jacobian} of~$f$. 



\subsection{Composition}






\section{Machine learning models as functions}

One way I understand many machine learning models, including deep
learning models, is that they are \emph{functions}.  The input is some
``state of the world,'' an element of the set, $\mathcal{X}$ of all possible
states, and the output is a ``prediction,'' an element of the set,
$\mathcal{Y}$, of all possible predictions.  \emph{Example:} An image
classification model takes, as input, an image, and outputs a
probability distribution over image labels.  \emph{Example:} An LLM
takes, as input, a sequence of tokens, and outputs a probability
distribution over tokens.

A function is a map from one space, $\mathcal{X}$, the domain, to another,
$\mathcal{Y}$, the codomain.  What kind of spaces are the domains and codomains
of machine learning models?  Consider a black-and-white image of size
$w\times h$ pixels, in which each pixel has an intensity between $0$
and~$1$.  Write $\set{I}$ for the unit interval, $\set{I}=[0,1]$.  An
image is a choice of an element of $\set{I}$ for each pixel, of which
there are $w\times h$; thus the domain
is~$\mathcal{X}=\set{I}^{w\times h}$.  On the output side, suppose there are
$l$ possible labels that the model might predict.  A probability
distribution over labels is a number between $0$ and $1$ for each
label, such that the sum of all the numbers is~$1$.  Thus the codomain
is the set, $\mathcal{Y}= \Delta^{l-1}$, the unit
$(l-1)$-simplex.\sidenote{See, e.g.,
  \url{https://en.wikipedia.org/wiki/Simplex}.}

Notice that neither $\set{I}^{w\times h}$ nor $\Delta^{l-1}$ is a vector
space.\sidenote{Neither is a vector space because there is no natural
  notion of addition or multiplication by a number. What would a pixel
  intensity of $2$ represent? Or multiplication by~$-1$?}
Nonetheless, it is a polite fiction to suppose that all machine
learning maps are from one real vector space to another.  One way this
is done is by embedding the actual domain or codomain in $\set{R}^n$,
for some~$n$.  For example, since $I$ is naturally a subset of
$\set{R}$, we can take $\set{I}^{w\times h}$ to be a subspace of
$\set{R}^{w\times h}$ and $\set{I}^l$ to be a subspace of~$\set{R}^l$.  One
also feels free to avail oneself of the vector space structure in
$\set{R}$ even though there was no such structure in~$\set{I}$.

One must take care that the final output of the model lies within the
true codomain.  For example, classification models, like the image
classifier, typically have as their penultimate output an element of
$y\in\set{R}^n$ whose individual components $y_i$ may take on arbitrary
real values.  In order to construct an element of $\Delta^{l-1}$, it is
usual to exponentiate each component (resulting in a positive number)
and then normalise the sum of all components to be~$1$.  That is, one
computes, for the $i$th component $y_i$, the value
$(\exp {y_i})/ \sum_i \exp {y_i}$.\sidenote{This transformation is
  called the ``softmax'' and the $y_i$s are called the ``unnormalised
  log probabilities.''}

Thus, a machine learning model is taken to be a function from
$\set{R}^M$ (for some $M$) to $\set{R}^N$ (for some $N$).

There is one more space of importance.  A particular machine learning
model is of course really a parameterised family of models and the set
of parameters, $\\Pi$,,is itself usually taken to be some real vector
space, say~$\Pi=\set{R}^P$ (for some natural number $P$).  Fixing a
particular $p\in\set{R}^P$ one obtains a particular instance of a model.

Thus, a machine learning model is a function,
$\mu\colon \Pi\times\mathcal{X}\to\mathcal{Y}$, or equivalently,
$\mu\colon \set{R}^P\times\set{R}^M\to \set{R}^N$.  For a given ``state of the
world'' $x\in\set{R}^M$, and parameters $\theta\in \set{R}^P$, one obtains a
``prediction,'' $y\in\set{R}^N$ by $y = \mu(\theta; x)$.

\section{Gradient descent}

The central problem in machine learning is this: Given
\begin{enumerate}
\item a parameterised
family of functions, $\mu : \Pi\times\mathcal{X}\to\mathcal{Y}$; and
\item a set of pairs, 
$(\tilde{x}_1, \tilde{y}_1)$, $(\tilde{x}_2, \tilde{y}_2)$, \dots,
  $(\tilde{x}_n, \tilde{y}_n)$, where $\tilde{x}_i\in\mathcal{X}$ and
  $\tilde{y}_i\in\mathcal{Y}$,
\end{enumerate}
choose a specific parameter $\theta\in\Pi$---that is, choose a particular
$\mu$---so that the $\mu(\theta; \tilde{x}_i)$ are ``close to''
the~$\tilde{y}_i$.  There are different versions of ``close to,'' but
one usually writes down a real-valued function
$L\colon \mathcal{Y}^n\times\set{\mathcal{Y}}^n\times\Pi\to\set{R}$ and attempts to minimise
\begin{equation*}
  \label{eq:loss}
  \mathcal{L}(\theta) = L\bigl(\tilde{y}_1,\dotsc,\tilde{y}_n; \mu(\theta,\tilde{x}_1), \dotsc,
  \mu(\theta,\tilde{x}_n); \theta\bigr),
\end{equation*}
with respect to~$\theta$.  (The explicit dependence on $\theta$ allows one to
specify a preference for some parameters over others.)  The minimum is
usually found by gradient descent, each step of which requires one to
evaluate $\Delta_\theta \mathcal{L}$, the gradient of $\mathcal{L}$ with respect to~$\theta$.





\end{document}
